\begin{frame}{Panels}
    \begin{panel}{A titled panel}
        \begin{itemize}
            \item \texttt{panel} takes a title as its mandatory argument.
            \item It breaks across frames, so a long list will not overflow.
            \item Any \texttt{tcolorbox} key goes in the optional argument.
        \end{itemize}
    \end{panel}
\end{frame}

\begin{frame}{Two panels on one slide}
    \begin{panel}[height=2.6cm]{Fixed height}
        Pass \texttt{height} to line up panels that sit next to each other.
    \end{panel}

    \vspace{0.3cm}

    \begin{panel}[height=2.6cm]{Second panel}
        Colours all come from \texttt{vars/colors.tex}; change one there and
        every panel follows.
    \end{panel}
\end{frame}

\begin{frame}{Split panel}
    \begin{panelsplit}
        \centering
        \tikz[baseline]\node[draw=gray!40, fill=surface, rounded corners=2pt,
            minimum width=3.2cm, minimum height=2.4cm, text=gray]{your figure here};
        \tcblower
        \begin{itemize}
            \item \texttt{panelsplit} splits left and right at \texttt{\textbackslash tcblower}.
            \item The left column takes a third of the width by default.
            \item Override it with \texttt{[lefthand ratio=0.5]}.
        \end{itemize}
    \end{panelsplit}
\end{frame}

\begin{frame}{A figure on its own}
    \begin{panel}{Figure}
        \centering
        \tikz[baseline]\node[draw=gray!40, fill=surface, rounded corners=2pt,
            minimum width=6cm, minimum height=3cm, text=gray]{your figure here};
    \end{panel}
\end{frame}
